% Two-column comparison table: OSNIP vs fastProve
% Compile: xelatex osnip_vs_fastprove_table.tex
\documentclass[11pt,a4paper]{article}
\usepackage[margin=2.0cm]{geometry}
\usepackage{fontspec}
\usepackage{xeCJK}
\setmainfont{Times New Roman}
\setCJKmainfont{PingFang SC}
\usepackage{booktabs,array,xcolor,colortbl,amsmath}
\usepackage{hyperref}

\definecolor{headerbg}{HTML}{F0F0F0}
\definecolor{note}{HTML}{555555}

\title{\textbf{OSNIP vs fastProve：精度损失对照表}\\
\large 聚焦 Llama-3.2-3B-Instruct}
\author{fastProve evaluation}
\date{\today}

\begin{document}
\maketitle
\vspace{-1.2em}

\noindent
{\small\color{note}
\textbf{说明：} 左列取自 OSNIP 论文（Cao et al.，Tables 1--3）；
右列取自本仓库大规模实测
\texttt{results/raw/llama\_compare\_large.json}
（1500 条真实场景 prompt $\times$ 3 master keys，BF16，RTX 5090）。
两侧指标定义不同，已分行写清，避免混读。
}

\vspace{1.0em}

\begin{center}
\renewcommand{\arraystretch}{1.35}
\begin{tabular}{@{}>{\raggedright\arraybackslash}p{3.4cm}
  >{\raggedright\arraybackslash}p{5.6cm}
  >{\raggedright\arraybackslash}p{5.6cm}@{}}
\toprule
\rowcolor{headerbg}
\textbf{对比项} &
\textbf{OSNIP（论文）} &
\textbf{fastProve（本工作）} \\
\midrule
方案定位
  & 客户端 embedding 混淆语义零空间注入
  & 全层增广状态协变混淆（structural / full） \\
\addlinespace
同模型
  & Llama-3.2-3B-Instruct
  & Llama-3.2-3B-Instruct \\
\addlinespace
主要精度定义
  & 相对非私有基线的\textbf{任务准确率保留}
  & 相对明文路径的\textbf{轨迹一致率} \\
\addlinespace
任务平均保留 RP
  & \textbf{98.49\%}
    （10 项任务 Avg Acc：$0.597\to0.588$）
  & structural \textbf{100.1\%}
    \newline full \textbf{97.46\%}
    （同 10 项 $n{=}200$，\texttt{osnip\_style\_bench\_n200.json}） \\
\addlinespace
MMLU（3B）
  & $0.606\to0.572$（约 $-3.4$\,pp）
  & plain $0.295$ / struct $0.295$ / full $0.315$
    （$n{=}200$ 子集，绝对分协议不同） \\
\addlinespace
与明文 token 一致率
  & ---
    （论文不以此为主指标）
  & structural \textbf{96.26\%}
    \newline full \textbf{89.39\%}
    （3 keys 平均） \\
\addlinespace
生成整段一致
  & 摘要 BERTScore RP $\approx$ \textbf{100\%}
    （相对非私有任务分，非轨迹对齐）
  & structural \textbf{84.51\%}
    \newline full \textbf{63.13\%}
    （与明文 greedy 后缀 exact match） \\
\addlinespace
PPL 相对增幅 $\Delta\%$
  & WikiText-2：\textbf{$+31\%$}
    （$9.63\to12.58$）
  & structural \textbf{$\approx 0\%$}
    \newline full \textbf{$+5.79\%$}
    （1500 prompts 上 mean） \\
\addlinespace
同论文其他基线（3B PPL）
  & Cape $+913\%$，DynText $+1066\%$，
    InferDPT $+29699\%$
  & 远低于上述量级 \\
\addlinespace
数据规模（精度侧）
  & 标准基准 + WikiText-2 + 1000 摘要样本等
  & 1500 真实场景 prompt $\times$ 3 keys \\
\bottomrule
\end{tabular}
\end{center}

\vspace{1.2em}
\subsection*{对照结论}
\begin{enumerate}
  \item \textbf{PPL：} 在 Llama-3.2-3B 上，fastProve（struct.\ $\approx0\%$，full $+5.8\%$）
    的困惑度膨胀\textbf{小于} OSNIP 论文的 $+31\%$。
  \item \textbf{任务准确率：} OSNIP 给出近无损任务 RP（$98.5\%$）；
    fastProve \textbf{尚无}同协议任务榜，不能据此声称全面更优。
  \item \textbf{轨迹对齐：} 仅 fastProve 直接报告与明文 token/greedy 一致率
    （struct.\ $\sim96\%/85\%$；full $\sim89\%/63\%$）。
\end{enumerate}

\vspace{0.5em}
{\scriptsize\color{note}
来源：OSNIP Tables 1--3（Llama-3.2-3B-Instruct）；
fastProve \texttt{llama\_compare\_large.json}（3 keys 平均）。
}
\end{document}
